\section{The joint factor-coordinate dynamics}
\label{sec:joint-dynamics}

We first isolate the exact algebra, independently of the prime weights.

\begin{proposition}[Primitive-ray and radial coordinates]
\label{prop:coordinate-bijection}
The map
\begin{equation}
 (m,n)\longmapsto
 \left(a=\frac{n}{(m,n)},\ b=\frac{m}{(m,n)},\ k=(m,n)\right)
 \label{eq:coordinate-map}
\end{equation}
is a bijection from $\N^2$ to the triples
\[
 (a,b,k)\in\N^3,\qquad (a,b)=1.
\]
Moreover,
\begin{equation}
 mn=abk^2,
 \qquad \frac nm=\frac ab,
 \qquad (m,n)=k.
 \label{eq:coordinate-identities}
\end{equation}
Consequently the joint frequency map
\begin{equation}
 (a,b,k)\longmapsto
 \lambda_{a,b,k}:=\bigl(\log(a/b),\log k\bigr)
 \label{eq:joint-frequency-map}
\end{equation}
is injective.
\end{proposition}

\begin{proof}
The GCD decomposition gives positive coprime $a,b$ and recovers
$n=ak$, $m=bk$, proving bijectivity and
\eqref{eq:coordinate-identities}.  If two joint frequencies agree, then
$a/b=c/d$.  Reduced positive fractions are unique, so $(a,b)=(c,d)$;
equality of $\log k$ and $\log\ell$ then gives $k=\ell$.
\end{proof}

The ratio-only map is the projection of \eqref{eq:joint-frequency-map}
onto its first coordinate.  Its fibre over a primitive ray $(a,b)$ is
the entire supported set of radial integers $k$.  This is exactly the
compression diagnosed in TPC-12.

\subsection{A commuting unitary representation}

Let $\mathscr I\subset\{(a,b,k):(a,b)=1\}$ be any finite index set and
put $\cH_{\mathscr I}=\ell^2(\mathscr I)$.  On its standard basis define
two diagonal self-adjoint operators
\begin{equation}
 L_{\rm ray}e_{a,b,k}=\log(a/b)e_{a,b,k},
 \qquad
 L_{\rm rad}e_{a,b,k}=\log k\,e_{a,b,k}.
 \label{eq:joint-generators}
\end{equation}
They commute, and hence generate the unitary $\R^2$-action
\begin{equation}
 U(t,s)=\exp\!\bigl(i tL_{\rm ray}+i sL_{\rm rad}\bigr).
 \label{eq:joint-unitary-action}
\end{equation}
For a coefficient vector $x=(x_{a,b,k})$, the coordinate-sum observation
is
\begin{equation}
 \mathscr O_x(t,s)
 :=\sum_{(a,b,k)\in\mathscr I}
 x_{a,b,k}e^{it\log(a/b)+is\log k}.
 \label{eq:coordinate-sum-observation}
\end{equation}
The arithmetic transform \eqref{eq:intro-joint-expansion} is this
observation with the finite set and coefficient vector selected by
$X,h,D,W$.

\begin{proposition}[Projection loss and joint completeness]
\label{prop:projection-and-completeness}
For a fixed primitive ray $(a,b)$, the restriction $s=0$ observes only
\begin{equation}
 e^{it\log(a/b)}\sum_kx_{a,b,k}.
 \label{eq:one-coordinate-fibre-sum}
\end{equation}
In contrast, the full function $(t,s)\mapsto\mathscr O_x(t,s)$ uniquely
determines every coefficient $x_{a,b,k}$.
\end{proposition}

\begin{proof}
The first assertion is immediate.  For the second, the exponentials have
distinct frequency vectors by \cref{prop:coordinate-bijection}.  A
finite linear combination of distinct characters of $\R^2$ cannot
vanish identically.  Explicit finite-resolution, Fej\'er-weighted
recovery with quantitative sector thresholds is proved in
\cref{thm:exact-sector-recovery}.
\end{proof}

\begin{remark}[What ``dynamics'' means here]
The action \eqref{eq:joint-unitary-action} is an exact diagonal unitary
action on the augmented factor-coordinate coefficient space.  Its generators and joint
eigenvalues are explicitly defined logarithmic coordinates.  They are
not asserted to be Koopman or transfer-operator eigenvalues of the
logistic map, and they have no claimed identification with Riemann-zero
ordinates.  This separation of meanings is essential to the scope of
the paper.
\end{remark}

\begin{remark}[Why a product phase is not the same lift]
Adding $(mn/X)^{is}$ would attach the frequency
$\log(abk^2/X)=\log(ab/X)+2\log k$.  The unknown ray-dependent term
$\log(ab/X)$ prevents a single global linear recombination with
$\log(a/b)$ from producing $\log k$.  Once the primitive ray is known,
one can nevertheless subtract that term and divide by two.  The GCD
phase is therefore a direct ray-adapted observable, not a claim of more
information than the complete ratio--product two-Mellin signal.
\end{remark}
